\chapter{Pronósticos y regresión}
\label{ch:pronosticos}

\section{Una previsión es una hipótesis operativa}
Los pronósticos alimentan inventarios, capacidad y planificación. Una recta que
ajusta bien el pasado no garantiza precisión futura
\cite{hyndman2021forecasting}. QSB implementa tendencia lineal, media móvil,
suavizamiento exponencial simple, descomposición estacional multiplicativa y
regresión lineal por mínimos cuadrados. El JSON es una solicitud: contiene
modelo, método, horizonte y parámetros. Así se conserva la elección del método
junto con los datos. Los backends nativo y de validación están disponibles;
la integración externa es una extensión futura.

\section{Datos y errores}
Las observaciones se procesan en el orden de sus filas; las etiquetas no
reordenan fechas ni completan períodos ausentes. Una serie mensual debe tener
espaciado mensual coherente. QSB devuelve valores ajustados, residuos
\(e_t=y_t-\hat y_t\), MAD, MSE y MAPE:
\[
 \mathrm{MAD}=\frac1m\sum|e_t|,\quad
 \mathrm{MSE}=\frac1m\sum e_t^2,\quad
 \mathrm{MAPE}=\frac{100}{m_+}\sum_{|y_t|>10^{-12}}|e_t/y_t|.
\]
\(m\) cuenta las filas ajustadas disponibles, no necesariamente toda la
historia; \(m_+\) excluye valores reales con magnitud no mayor que \(10^{-12}\).
Si todos quedan excluidos, MAPE no
está disponible. MAD conserva unidades de la serie; MSE tiene unidades al
cuadrado; MAPE se publica en porcentaje. Comparar métodos exige considerar
que media móvil y suavizamiento omiten filas iniciales diferentes.

\section{Tendencia lineal: ejemplo completo}
El modelo \(y_t=a+bt+e_t\) estima \(a,b\) por mínimos cuadrados con índices
\(t=1,\ldots,n\). \path{manual/examples/prob_pronostico.json} contiene
\((12,14,16,18,20,22)\), unidad mes y horizonte dos. Desde la raíz:
\begin{lstlisting}
.build/debug/qsb validate-forecast-json manual/examples/prob_pronostico.json
.build/debug/qsb solve-forecast-json manual/examples/prob_pronostico.json
\end{lstlisting}
QSB obtiene \(a=10,b=2\), pronósticos \(24\) y \(26\), y los tres errores
cero. Son datos artificiales elegidos para comprobar fórmula e índices: el
error nulo no constituye evidencia sobre demanda real. El resultado conserva
la solicitud, el algoritmo \texttt{linearTrend} y exactitud
\texttt{exact}, que describe el cálculo del método, no el conocimiento exacto
del futuro.

\section{Media móvil y suavizamiento}
Para ventana \(k\), la previsión de un paso es
\[
 \hat y_{t+1}=\frac1k\sum_{j=0}^{k-1}y_{t-j}.
\]
QSB extiende horizontes usando previsiones anteriores dentro de la ventana
cuando se agotan valores reales. El archivo \path{prob_movil.json}
selecciona ventana tres sobre la misma serie. Resultan \(20\) y
\(20.666667\): la segunda media utiliza \((20,22,20)\). MAD es \(4\), MSE
\(16\) y MAPE \(20.134680\%\), calculados sobre las últimas tres filas.
Esta demora es previsible cuando la serie tiene tendencia.

El suavizamiento simple aplica
\[
 F_{t+1}=\alpha y_t+(1-\alpha)F_t,\qquad 0<\alpha\le1.
\]
El pronóstico inicial se fija al primer dato. En
\path{prob_suavizado.json}, \(\alpha=0.3\) produce \(18.117660\) en
ambos períodos futuros. El método no tiene componente de tendencia; no
extrapola automáticamente el crecimiento. Su MAD es \(4.078440\), MSE
\(18.233818\) y MAPE \(22.081810\%\). Un alfa grande reacciona rápido y
puede seguir más ruido; uno pequeño suaviza y puede retrasarse.
Ambos ejemplos se ejecutan con \texttt{solve-forecast-json} y se validan con
\texttt{validate-forecast-json}.

\section{Estacionalidad multiplicativa}
El método combina la tendencia estimada con factores de estación:
\(\hat y_t=(a+bt)s_j\). QSB calcula cocientes real/tendencia, promedia por
posición estacional y normaliza los factores para que su media sea uno.
Es una descomposición educativa, no Holt--Winters ni un modelo ARIMA. La
longitud estacional y el orden de observaciones deben corresponder al fenómeno;
se necesitan al menos dos ciclos completos y valores/tendencia compatibles con
la división multiplicativa.
\path{prob_estacional.json} usa longitud dos sobre la serie lineal anterior.
Los dos factores son uno; la previsión vuelve a ser \(24,26\). El caso
comprueba que especificar una estación no inventa oscilaciones cuando los datos
coinciden con la recta. Para estudiar estacionalidad real, cambie la serie por
varios ciclos con oscilación y revise residuos antes de interpretar los factores.

\section{Regresión explicativa}
Se modela \(y_i=\beta_0+\sum_j\beta_jx_{ij}+e_i\) con intercepto.
Los predictores deben tener columnas linealmente independientes; una columna
constante duplica el intercepto. El núcleo detecta deficiencia de rango.
\path{prob_regresion.json} contiene actividad \(1,2,3,4,5\) y consumo
\(5,7,9,11,13\). Con \texttt{solve-forecast-json}, QSB devuelve intercepto
\(3\), coeficiente de actividad \(2\), SSE cero y \(R^2=1\).
El horizonte del contrato de solicitud debe ser positivo también para
regresión, aunque su salida ajusta las observaciones suministradas y no genera
filas futuras de predictores. Para una nueva actividad, la ecuación estimada
permite calcular una predicción, sin que QSB publique un intervalo para ella.

Un coeficiente es una asociación condicional, no causalidad. \(R^2\) mide
ajuste dentro de la muestra; no identifica sobreajuste ni garantiza
extrapolación. No se publican errores estándar de coeficientes, pruebas de
hipótesis ni intervalos predictivos en este flujo. Datos originales pequeños
permiten comprobar el contrato, pero no evaluar capacidad de generalización.

\section{Trabajo en macOS y evaluación}
En el editor de pronósticos seleccione método, parámetros y horizonte; introduzca
filas históricas en el orden temporal. Se pueden pegar valores en una columna,
o etiquetas y valores separados por tabuladores en dos columnas. El pegado se
valida de forma atómica: una fila mal formada no reemplaza parcialmente la
historia. La regresión utiliza nombres de variables y filas explicativas.
Valide, resuelva y revise las curvas reales, ajustadas y futuras; las curvas y
el gráfico de residuos comparten el orden de observaciones. Controles de
visibilidad y escala horizontal ayudan con series extensas. El resumen muestra
los parámetros y métricas; el selector JSON conserva solicitud y resultado.
Aplicar JSON solo actualiza el borrador cuando su decodificación funciona.

Para evaluar un pronóstico operacional, reserve períodos finales, ajuste con
la historia anterior y compare con datos reservados. Los errores publicados por
QSB son los ajustes del cálculo ejecutado; no constituyen automáticamente una
prueba fuera de muestra. No seleccione parámetros mirando reiteradamente el
mismo conjunto reservado. Documente unidades, ventana, alfa, estación y horizonte
al entregar el análisis.

\section{Ejercicios}
\begin{enumerate}
\item Compare ventanas dos y cuatro sobre la serie creciente. Explique el
retraso y compruebe qué filas participan en MAD.
\item Sustituya un dato por cero y compruebe la definición de MAPE. Compare
su conveniencia con MAD.
\item Introduzca un segundo predictor igual al doble del primero y observe el
diagnóstico de rango. Explique por qué más columnas no siempre aportan información.
\end{enumerate}
